\documentclass[border=8pt]{standalone}

% --- packages (mirror these in template.meta.json "requires") ---
\usepackage[dvipsnames]{xcolor}
\usepackage{tikz}
\usetikzlibrary{calc, backgrounds, fit, arrows.meta}

% --- palette: ESL/tsarilp domain-semantic scheme (NOT opentikz's ot* palette) ---
% Same role/color assignments as esl-architecture-block (core-accent=Aquamarine,
% inter-bus=NavyBlue dashed, intra-bus=Orange) so a \foreach-generated array reads
% as the same family, not a different visual language.

\begin{document}
% ==== parameters (edit these; this is the \foreach generator's parameter surface —
% nothing below this block should ever be edited to change array shape/size) =====
\def\ntilerows{2}       % number of tile rows in the MPSoC grid
\def\ntilecols{2}       % number of tile columns in the MPSoC grid
\def\pegridrows{2}      % PEs per tile, arranged as a \pegridrows x \pegridcols sub-grid
\def\pegridcols{2}
\def\tilelabelprefix{T} % tile label prefix; rendered as "<prefix><row>.<col>"
\def\pelabelmode{index} % index|none -- legibility control for PE-node text at scale
\def\tilegapx{0.9}      % gap between tile boundary boxes, horizontal (cm)
\def\tilegapy{0.9}      % gap between tile boundary boxes, vertical (cm)
\def\pegap{0.15}        % gap between PEs within a tile (cm)
\def\pesize{0.5}        % PE node minimum size (cm)
\def\showmesh{1}        % 1 = draw inter-tile mesh bus (NoC) to right/below neighbor; 0 = omit
% =========================================================================

% legibility sentinel: \pelabelmode is compared against this literal, so a PE node's
% label is drawn only when the author opted into "index" mode (see invariants).
\def\PEShowIndex{index}

\begin{tikzpicture}[
    >=Stealth,
    pe/.style={draw, thick, rectangle, minimum size=\pesize cm, inner sep=0pt,
      draw=Aquamarine!80!black, fill=Aquamarine!25, font=\tiny, align=center},
    tilebox/.style={draw=Aquamarine!70!black, dashed, rounded corners=3pt,
      inner sep=6pt},
    tilelabel/.style={font=\scriptsize\bfseries, text=Aquamarine!55!black},
    intrabus/.style={draw, Orange, thick},
    interbus/.style={draw, NavyBlue, thick, dashed},
  ]

  % --- pitch: uniform spacing between successive tile origins, derived from the
  % PE sub-grid extent plus the inter-tile gap. Nothing above this arithmetic
  % reasons in raw (x,y); the author only ever states counts/gaps above. ---
  \pgfmathsetmacro{\pestepx}{\pesize+\pegap}
  \pgfmathsetmacro{\pestepy}{\pesize+\pegap}
  \pgfmathsetmacro{\tilepitchx}{\pegridcols*\pestepx+\tilegapx}
  \pgfmathsetmacro{\tilepitchy}{\pegridrows*\pestepy+\tilegapy}

  % ==== generator: tile grid x { PE sub-grid, containment fit box } ==========
  \foreach \r in {1,...,\ntilerows} {
    \pgfmathsetmacro{\oy}{-(\r-1)*\tilepitchy}
    \foreach \c in {1,...,\ntilecols} {
      \pgfmathsetmacro{\ox}{(\c-1)*\tilepitchx}
      \def\tilefitlist{}
      \foreach \pr in {1,...,\pegridrows} {
        \foreach \pc in {1,...,\pegridcols} {
          \pgfmathtruncatemacro{\k}{(\pr-1)*\pegridcols+\pc}
          \pgfmathsetmacro{\px}{\ox+(\pc-1)*\pestepx}
          \pgfmathsetmacro{\py}{\oy-(\pr-1)*\pestepy}
          \node[pe] (pe-\r-\c-\k) at (\px,\py) {%
            \ifx\pelabelmode\PEShowIndex \k\fi
          };
          \xdef\tilefitlist{\tilefitlist(pe-\r-\c-\k)}
        }
      }
      % containment: the fit box is declared after every PE it fits, on the
      % background layer, so it sits behind them (primitive: containment).
      \begin{scope}[on background layer]
        \node[tilebox, fit=\tilefitlist] (tile-\r-\c) {};
      \end{scope}
      \node[tilelabel, anchor=north west]
        at ([xshift=2pt,yshift=-2pt]tile-\r-\c.north west)
        {\tilelabelprefix\r.\c};
      % intra-tile bus: a light reading-order chain across this tile's PEs
      % (primitive: relative-placement / typed-routing within the container).
      \pgfmathtruncatemacro{\npes}{\pegridrows*\pegridcols}
      \ifnum\npes>1
        \foreach \k in {2,...,\npes} {
          \pgfmathtruncatemacro{\kprev}{\k-1}
          \draw[intrabus] (pe-\r-\c-\kprev) -- (pe-\r-\c-\k);
        }
      \fi
    }
  }

  % ==== inter-tile mesh (NoC): each tile wires to its right/below neighbor only
  % -- O(tiles) edges, not O(tiles^2), so the mesh stays legible at scale
  % (primitive: alignment-symmetry / typed-routing across containers) ====
  \ifnum\showmesh=1
    \foreach \r in {1,...,\ntilerows} {
      \foreach \c in {1,...,\ntilecols} {
        \ifnum\c<\ntilecols
          \pgfmathtruncatemacro{\cnext}{\c+1}
          \draw[interbus] (tile-\r-\c.east) -- (tile-\r-\cnext.west);
        \fi
        \ifnum\r<\ntilerows
          \pgfmathtruncatemacro{\rnext}{\r+1}
          \draw[interbus] (tile-\r-\c.south) -- (tile-\rnext-\c.north);
        \fi
      }
    }
  \fi

\end{tikzpicture}
\end{document}
